\subsection{Higher Hamiltonians and one-point BA times}
\label{sec:BA-higher-flows}

The commuting Hamiltonians of \cref{sec:spectral-flows} act on
the same BA data in different directions. We now compute those
directions and the accompanying constant mode phases. Jacobian
linearization and its description by principal parts are classical
\cite[\S\S6--7]{Griffiths1985}; see also
\cite{Beauville1990,AdamsHarnadHurtubise1993}.
Here a gap normalization converts the direction into a one-point BA
time, and finite parts give the explicit phases of the canonical
modes. The construction applies to every prescribed realization,
including coincident phases and singular plane spectral images.

Fix a complex spectral value $\nu_0$ and an integer $r\geq2$. On the
$H_0$ pencil \cref{eq:BA-H0-pencil}, put
\begin{equation}\label{eq:BA-shifted-flow}
\begin{aligned}
 H_{r,\nu_0}&=\tfrac1r\tr(\mathsf X+\nu_0\mathsf D)^r,
 &\mathsf G&=(\mathsf X+\nu_0\mathsf D)^{r-1},\\
 \mathsf T_{r,\nu_0}(\nu)
 &=\frac{\mathsf G-\mathsf L_0(\nu)^{r-1}}{\nu-\nu_0},
 &q(P)&=\frac{\lambda(P)^{r-1}}{\nu(P)-\nu_0}.
\end{aligned}
\end{equation}
The quotient defining $\mathsf T_{r,\nu_0}$ is polynomial, and
$i\dot{\mathsf L}_0=[\mathsf L_0,\mathsf T_{r,\nu_0}]$.
Indeed $[\mathsf L_0,\mathsf G]=(\nu-\nu_0)[\mathsf D,\mathsf G]$.
All coefficients of $H_{r,\nu_0}$ in $\nu_0$ are commuting spectral
Hamiltonians. The unshifted traces are $H_{r,0}$.

For $g\geq1$, let $n_1<\cdots<n_g\leq2g-1$ be the Weierstrass
gaps at $P_\infty$: these are the positive integers that cannot be
the pole order of a function having its only pole there. Write
$\operatorname{pp}_Q$ for the principal part at a point $Q$.

\begin{lemma}[One-point normalization]\label{lem:BA-higher-normalization}
There is a unique meromorphic function $\mathfrak f$ with poles only
at $Q_1,\ldots,Q_K,P_\infty$ such that
\[
 \operatorname{pp}_{Q_k}\mathfrak f=-\operatorname{pp}_{Q_k}q,
 \quad1\leq k\leq K,\qquad
 \mathfrak f=-\sum_{\alpha=1}^{g}v_\alpha s^{-n_\alpha}+O(s)
 \quad(P\to P_\infty).
\]
Both $\mathfrak f$ and the coefficients $v_\alpha$ are obtained by a
finite linear meromorphic interpolation problem. In genus zero the
central condition is $\mathfrak f=O(s)$ and the sum is empty.
\end{lemma}
\begin{proof}
For $g\geq1$, Riemann--Roch gives
$H^1(\mathcal O((2g-1)P_\infty))=0$. Thus arbitrary outer principal
parts extend to a meromorphic function with an additional central
pole of order at most $2g-1$ \cite[Chapter~IV]{Hartshorne1977}.
Subtract functions with their only pole at $P_\infty$ to remove the
nongap powers in descending order, then subtract the central constant.
A difference of two normalized solutions would have its only pole
at $P_\infty$ with largest order a gap, which is impossible; its zero
constant then makes it vanish. On the sphere the outer principal
parts extend without a central pole, and fixing the central value
gives the same conclusion.
\end{proof}

Use $1/\nu$ as coordinate at the outer points. Put $d_i=1$ for
$1\leq i\leq L$ and $d_j=-1$ for $L+1\leq j\leq K$, and define
\begin{equation}\label{eq:BA-higher-time-phase}
 p(s^{-1})=\sum_{\alpha=1}^{g}v_\alpha s^{-n_\alpha},\qquad
 b_k^{(r,\nu_0)}=-d_k(\mathfrak f+q)(Q_k),\qquad
 c_k(t)=c_k e^{-it b_k^{(r,\nu_0)}},\quad1\leq k\leq K.
\end{equation}
The sum $\mathfrak f+q$ is regular at every $Q_k$. Its values and
the coefficients of $p$ depend only on the fixed marked spectral data.
Let $\boldsymbol U_{n_\alpha}$ be the BA direction of the normalized
second-kind differential with principal part $d(s^{-n_\alpha})$.
The split into $p$ and mode phases uses the fixed local parameter $s$.
Changing that parameter changes this representation of the same
Hamiltonian evolution.

\begin{theorem}[The spectral BA hierarchy]\label{thm:BA-higher-flows}
Every prescribed realization in \cref{thm:general-BA-realization}
evolves under $H_{r,\nu_0}$ by the same ordinary and dual BA functions
and Cauchy kernel, with one-point time $p$ and gauges $c_k(t)$ from
\cref{eq:BA-higher-time-phase}. Explicitly, use those gauges in
\cref{eq:general-elementary-modes,eq:general-pair-modes}, and replace
the first-time exponentials by $e^{-itp}$ and $e^{itp}$ at the puncture.
The Abel argument is
\begin{equation}\label{eq:BA-higher-Abel-direction}
 z(t)=z_0-it\sum_{\alpha=1}^{g}v_\alpha\boldsymbol U_{n_\alpha}.
\end{equation}
The elementary partner phases are reciprocal; each pair mode has
phase $e^{it(b_i^{(r,\nu_0)}+b_j^{(r,\nu_0)})}$ and its partner
has the inverse phase. The statement holds locally on the nonspecial
BA chart, without a distinct-phase assumption.
Finite constant linear combinations of the Hamiltonians use the
corresponding sums of $p$ and the phase rates.

As $\nu_0$ varies, the flows with $2\leq r\leq K+1$ span all
Jacobian directions. Consequently finitely many commuting polynomial
spectral coefficients suffice to generate those translations.
\end{theorem}
\begin{proof}
Start with the prescribed BA modes and the right eigenvector
\cref{eq:BA-prescribed-eigenvector}. The left vector is
\[
 \ell(P)=\bigl((-c_i\Omega(P,Q_i))_{i=1}^{L},
 (\kappa_j\Omega(P,Q_j)/c_j)_{j=L+1}^{K},\psi(P)\bigr).
\]
The convolution identities give its spectral equation and
\begin{equation}\label{eq:BA-higher-initial-kernel}
 \ell(P)r(Q)=-(\nu(P)-\nu(Q))\Omega(P,Q).
\end{equation}
Evolve the matrix by the canonical Hamiltonian and these vectors by
$i\dot r_H=-\mathsf T_{r,\nu_0}r_H$ and
$i\dot\ell_H=\ell_H\mathsf T_{r,\nu_0}$. Polynomiality of the time
matrix preserves their finite pole divisors.

At infinity, write $\partial=i\partial_t$. The eigenvalue equation
gives the convergent local factorization
\begin{equation}\label{eq:BA-higher-local-factorization}
 r_H=e^{-itq(P)}\mathsf V(t,\nu(P))r(P,0),\qquad
 \partial\mathsf V=-\frac{\mathsf G}{\nu-\nu_0}\mathsf V,
 \qquad \mathsf V(0,\nu)=\Id.
\end{equation}
Since $\mathsf V=\Id+O(\nu^{-1})$, every leading coordinate direction
is preserved, even when branch intercepts coincide. Dress the vectors
reciprocally by
$\widetilde r=e^{-it\mathfrak f}r_H$ and
$\widetilde\ell=e^{it\mathfrak f}\ell_H$.
At the central point $q=O(\nu^{-1})$, so the last components have
the normalized dual and ordinary factors $e^{itp}$ and $e^{-itp}$;
all other components are $O(s)$ times those factors.
At each outer point the polar exponent cancels. If $R_k$ is the
finite part of $q+\mathfrak f$, the right vector's leading
coefficients are
\[
 -\nu e^{-itR_i}/c_i\quad(1\leq i\leq L),\qquad
 c_j\nu e^{-itR_j}/\kappa_j\quad(L+1\leq j\leq K).
\]
These are precisely $-\nu/c_i(t)$ and $c_j(t)\nu/\kappa_j$.

BA uniqueness identifies the last components. Each remaining component
has the fixed BA divisor, one additional simple pole at its own
outer point, the calculated residue and zero leading coefficient at
the central puncture. The same uniqueness identifies the kernel
components. No new poles arise at distinct points of equal $\nu$:
the derivative of $\ell_H(P)r_H(Q)$ is divisible by
$\nu(P)-\nu(Q)$, because the time matrix is polynomial. Thus all
initially removable poles in \cref{eq:BA-higher-initial-kernel}
remain removable. The diagonal pairing is constant, so the duality
differential and weights are unchanged. The central row and the
opposite outer blocks of the spectral equation now recover all modes
as in \cref{eq:BA-converse-recovery}. This proves their asserted
BA formulas, including every pair phase.

For the spanning assertion, choose an unramified finite fiber
$\nu=\nu_0$ with distinct eigenvalues $\lambda_h$, $1\leq h\leq K+1$.
The function $\mathfrak f+q$ has poles only at the central point,
with principal part $-p$, and on this fiber, with the principal parts
of $q$. For a holomorphic differential $du=w(P)d\nu$, the residue
theorem therefore gives
\[
 \operatorname{Res}_{P_\infty}(p\,du)
 =\sum_{h=1}^{K+1}\lambda_h^{r-1}w(P_h).
\]
If $du$ annihilates every such BA direction, the left side vanishes
for $2\leq r\leq K+1$. The same vanishing for $r=1$ follows by
applying the residue theorem to $du/(\nu-\nu_0)$.
The Vandermonde matrix is invertible, so $w$ vanishes on every such
fiber, hence $du=0$. Residue duality proves spanning.
Finite-dimensionality and expansion in $\nu_0$ give the final assertion.
\end{proof}

On a fixed spectral level the vectors in
\cref{eq:BA-higher-Abel-direction} are constant, and constant linear
combinations of the Hamiltonians select any Abel direction.
The required combination depends on the spectral data: the theorem
identifies polynomial Hamiltonian flows with level-dependent
combinations of normalized BA times.

\begin{corollary}[Higher flows from generic initial modes]
\label{cor:BA-higher-reconstruction}
On the generic initial-data domains of \cref{thm:BA-spectral-converse},
the same reconstructed curve and BA data solve the commuting flows
$H_{r,\nu_0}$ with \cref{eq:BA-higher-time-phase}.
For repeated phases, use the same constant active bases and restore
the original coordinates after the BA evolution.
\end{corollary}
\begin{proof}
Apply \cref{thm:BA-higher-flows} to the reconstructed data. Every
spectral Hamiltonian is invariant under equal-phase block conjugations,
so their moment matrices are conserved and the diagonalizing bases
remain constant. Inactive directions are direct summands of the
pencil; matrix powers preserve their zero couplings. Thus the reduction
and restoration in \cref{eq:BA-block-bases,eq:BA-inactive-splitting}
apply to all these flows.
\end{proof}

For the fourth trace, $q=\lambda^3/\nu$. Its finite part has a useful
modal-power interpretation. Write $e_k=\rho-b_k$ and let
$C_i=I_i-\sum_{j=L+1}^{K}I_{ij}$,
$C_j=I_j-\sum_{i=1}^{L}I_{ij}$, where
$I_k=\xi_k\bar\xi_k$ and $I_{ij}=\xi_{ij}\bar\xi_{ij}$.
On a prescribed realization, the eigenvector's pure coordinate pole
at each $Q_k$ makes the corresponding moment block diagonal in that
basis, also when phases coincide.
Second-order perturbation at infinity gives
\[
 \lambda=d_k\nu+e_k-d_kC_k/\nu+O(\nu^{-2}),\qquad
 q=d_k\nu^2+3e_k\nu+3d_k(e_k^2-C_k)+O(\nu^{-1}).
\]
Thus the phase in \cref{eq:BA-higher-time-phase} is computed from
the outer finite part of $\mathfrak f$ and conserved power charges.
On a fixed spectral level these constant phase rates can also be
removed by subtracting $\sum_{k=1}^{K}b_k^{(r,\nu_0)}C_k$ from the
Hamiltonian. The resulting flow has the normalized one-point BA time.
